\section{Metatheory}\label{sec:meta}

In this section, we provide the denotational semantics of types
($\denot{\tau}$) and type contexts ($\denot{\Gamma}$), and state the
soundness theorem of our type system.  The key challenge is to give a
semantics that supports both demonic and angelic modalities
simultaneously, together with the resource-based interpretation of
contexts that tracks variable dependencies.  We proceed in four
steps: we define contextual capabilities as the semantic domain
(\S\ref{subsec:capabilities}), equip them with algebraic structure
(\S\ref{subsec:algebra}), build a context logic over this algebra
(\S\ref{subsec:logic}), and finally give the type denotation and
soundness theorem (\S\ref{subsec:denotation}).

\subsection{Contextual Capabilities}\label{subsec:capabilities}

In a pure functional language, the free variables of a term $e$ are
bound by its evaluation environment, a partial map from variables to
values.  To reason about both safety and reachability, we need to
consider not just a single environment but a \emph{set} of
environments, representing the range of possible variable bindings
available to demonic or angelic choices.  We call such a set a
\emph{contextual capability}.

\begin{definition}[Contextual Capability]\label{def:contextual-capability}
  A \emph{contextual capability} $r$ is a non-empty set of
  environments (partial maps $\Var \rightharpoonup \Val$) that all
  share the same domain.  We write $\Res$ for the set of all
  well-formed contextual capabilities and assume throughout that all
  capabilities are well-formed.
\end{definition}

The domain condition ensures that all environments in a capability
bind exactly the same set of variables, reflecting the fact that the
free variables of a term are fixed by its typing context.  A
capability with a single environment corresponds to a fully determined
variable binding; a capability with multiple environments represents
the range of choices available for demonic or angelic reasoning.

We extend the standard notions of environment compatibility
($r_1 \uparrow r_2$, meaning every environment in $r_1$ is compatible
with every environment in $r_2$), projection ($\projectTo{r}{X}$, the
piecewise restriction of $r$ to variable set $X$), and refinement
($r_1 \sqsubseteq r_2$, meaning $r_1 = \projectTo{r_2}{\Dom{r_1}}$)
from environments to capabilities in the obvious pointwise way.

The refinement relation $\sqsubseteq$ plays a central role: $r_1
\sqsubseteq r_2$ says that $r_1$ is a restriction of $r_2$ to a
smaller variable domain, i.e., $r_2$ extends the bindings of $r_1$
with additional variables.  This relation will serve as the Kripke
accessibility relation in our logic (\S\ref{subsec:logic}).

To express variable dependency in refinement type qualifiers, we also
need the following notion.

\begin{definition}[Fiber]\label{def:fibers}
  For a variable $x$ and an environment $[x \mapsto v_x] \in
  \projectTo{r}{\{x\}}$, the \emph{fiber} of $r$ over $[x \mapsto
  v_x]$ is:
  \[
    \fiberOver{r}{[x \mapsto v_x]} \doteq \{\, \sigma' \in r \mid
    \sigma'(x) = v_x \,\}.
  \]
  That is, $\fiberOver{r}{[x \mapsto v_x]}$ is the sub-capability of
  $r$ consisting of all environments that assign $x$ to $v_x$.
\end{definition}

\noindent Fibers allow us to condition a capability on the value of a
specific variable, which is essential for evaluating refinement
qualifiers that refer to other variables bound in the typing context
(see \S\ref{subsec:logic}).

\subsection{Capability Algebra}\label{subsec:algebra}

Our context logic is a bunched implication logic (BI-logic) based on a
BI-algebra~\cite{...}, a partial commutative monoid with a pre-order
satisfying bifunctoriality.  We equip contextual capabilities with two
operations: a \emph{product} that composes disjoint variable bindings,
and a \emph{sum} that merges capabilities over the same domain to
support control-flow reasoning.

We equip $\Res$ with two operations. The \emph{product} $r_1 \times
r_2$, defined when $r_1 \uparrow r_2$ is the set of all pairwise
unions $\sigma_1 \cup \sigma_2$ $\sigma_1 \in r_1$ and $\sigma_2 \in
r_2$; it combines capabilities over disjoint variable domains. The
\emph{sum} $r_1 + r_2$​, defined when $\Dom{r_1} = \Dom{r_2}$​, is
simply $r_1 \cup r_2$​; it merges capabilities over the same domain to
support control-flow reasoning. Together with unit $\mathbf{1} \doteq
\{\emptyset\}$, the tuple $(\Res, +, \times, \mathbf{1}) (\Res,+,×,1)$
forms a partial commutative semiring without zero. To enable a Kripke
semantics for BI-logic, this semiring must be equipped with a partial
order satisfying bifunctoriality.

The natural candidates for this order are the subset ($\subseteq$) and
superset ($\supseteq$) relations.  However, neither works for both
modalities simultaneously: the subset order forces an
overapproximation semantics (demonic), while the superset order forces
underapproximation (angelic).  Since we need to support both as
first-class operators, the Kripke order must be \emph{neutral} with
respect to approximation.  The refinement relation $\sqsubseteq$ has
exactly this property: $r1 \sqsubseteq r_2$ holds if $r_2$ extends
$r_1$ with additional bindings, while preserving all existing
environments, neither adding nor removing possible values for the
variables already in $r_1$'s domain.  Specifically, $\sqsubseteq$
satisfies the standard bifunctoriality property required of a
BI-algebra: $r_1 \sqsubseteq r_1'$ and $r_2 \sqsubseteq r_2'$ implies
$r_1 \times r_2 \sqsubseteq r_1' \times r_2'$, ensuring that $(\Res,
+, \times, \mathbf{1}, \sqsubseteq)$ forms a BI-algebra.


\subsection{Context Logic}\label{subsec:logic}

We now define a logic over the capability algebra.  The logic is a
bunched implication logic extended with demonic and angelic
modalities, a sum operator for control-flow reasoning, and a
\emph{binding reference operator} that expresses how a qualifier may
depend on specific variable bindings in the current capability.

\begin{definition}[Context Logic]\label{def:context-logic}
  The syntax of context logic is:
  \begin{align*}
    P, Q \;::=\; &\chtrue \mid \chfalse \mid P \chland Q \mid P \chlor Q \mid P \chimpl Q \mid P \chstar Q \mid P \chwand Q
    \\& \mid \forall x.\, P \mid \exists x.\, P \mid \Code{Atom}(\phi)
    \mid \demonic P \mid \angelic P \mid P \choplus Q \mid \forallM{x} P
  \end{align*}
  The forcing semantics over capabilities $r \in \Res$ is:
  \begin{itemize}
  \item $r \models \chtrue$ always; \quad $r \models \chfalse$ never;
  \item $r \models P \chland Q$ iff $r \models P$ and $r \models Q$;
  \item $r \models P \chlor Q$ iff $r \models P$ or $r \models Q$;
  \item $r \models P \chimpl Q$ iff for all $r \sqsubpseteq r'$, $r' \models P$ implies $r' \models Q$;
  \item $r \models P \chstar Q$ iff $\exists r_1, r_2.\; r_1 \times r_2 \sqsubseteq r$,\; $r_1 \models P$,\; $r_2 \models Q$;
  \item $r \models P \chwand Q$ iff $\forall r' \uparrow r$,\; $r' \models P$ implies $r' \times r \models Q$;
  \item $r \models \forall x.\, P$ iff $x \notin \Dom{r}$ and $\forall r \sqsubseteq r'$ with $x \in \Dom{r'}$,\; $r' \models P$;
  \item $r \models \exists x.\, P$ iff $x \notin \Dom{r}$ and $\exists r \sqsubseteq r'$ with $x \in \Dom{r'}$,\; $r' \models P$;
  \item $r \models \Code{Atom}(\phi)$ iff $\projectTo{r}{\Code{FV}(\phi)} = \{\,\sigma \mid \Dom{\sigma} = \Code{FV}(\phi) \land \sigma \models \phi\,\}$;
  \item $r \models \demonic P$ iff $\exists r' \supseteq r$ such that $r' \models P$;
  \item $r \models \angelic P$ iff $\exists r' \subseteq r$ such that $r' \models P$;
  \item $r \models P \choplus Q$ iff $\exists r_1, r_2.\; r_1 + r_2 \sqsubseteq r$,\; $r_1 \models P$,\; $r_2 \models Q$;
  \item $r \models \forallM{x} P$ iff $\forall [x \mapsto v_x] \in \projectTo{r}{\{x\}}.\; \fiberOver{r}{[x \mapsto v_x]} \models P[x \mapsto v_x]$.
  \end{itemize}
  We write $P \models Q$ iff for all $r \in \Res$, $r \models P$ implies $r \models Q$.
\end{definition}

The additive connectives ($\chland$, $\chlor$, $\chimpl$, $\forall$,
$\exists$) are standard.  The multiplicative connectives ($\chstar$,
$\chwand$) follow standard BI-logic semantics using the capability
product.  We discuss the remaining operators below.

\paragraph{Kripke Monotonicity}
The implication $P \chimpl Q$ and the first-order quantifiers
$\forall x.\, P$ and $\exists x.\, P$
use the refinement
relation $\sqsubseteq$ as the Kripke accessibility relation: truth at
$r$ requires truth at all $r \sqsubseteq r'$, i.e., at all extensions
of $r$ with additional variable bindings.  This is the standard Kripke
semantics for implication in BI-logic~\cite{...}, and ensures the
following monotonicity property.

\begin{theorem}[Kripke Monotonicity]\label{thm:kripke-monotonicity}
  $r \sqsubseteq r'$ and $r \models P$ implies $r' \models P$.
\end{theorem}

The use of $\sqsubseteq$ rather than $\subseteq$ or $\supseteq$ as the
accessibility relation is what allows our logic to remain neutral with
respect to approximation direction, as argued in
\S\ref{subsec:algebra}.

\paragraph{Atomic Propositions}
The atomic proposition $\Code{Atom}(\phi)$ lifts a standard Boolean
qualifier $\phi$ to a capability-level predicate with \emph{exact}
semantics: a capability $r$ satisfies $\Code{Atom}(\phi)$ iff its
projection to $\Code{FV}(\phi)$ contains \emph{exactly} those
environments that satisfy $\phi$.  This neutrality leaves the choice
of approximation direction entirely to the demonic and angelic
modalities.

\paragraph{Demonic and Angelic Modalities}
The demonic modality $\demonic P$ holds at $r$ if $r$ can be
\emph{extended} to some $r'$ (i.e., $r \sqsubseteq r'$) that satisfies
$P$; this is an overapproximation, since the extension may include
additional environments beyond what $r$ contains.  Dually, the angelic
modality $\angelic P$ holds at $r$ if some \emph{sub-capability} $r'
\sqsubseteq r$ satisfies $P$; this is an underapproximation, since $r'$
witnesses a specific subset of the environments in $r$.

\begin{example}[Approximation Modalities]\label{ex-approximation-modalities}
  Let $r_x = \{[x \mapsto 1], [x \mapsto 2]\}$.  The following hold:
  \begin{align*}
    &r_x \models \demonic\,\Code{Atom}(x = 1) \;\judgfail \quad
    &&r_x \models \demonic\,\Code{Atom}(0 \leq x \leq 4) \;\judgok \\
    &r_x \models \angelic\,\Code{Atom}(x = 1) \;\judgok \quad
    &&r_x \models \angelic\,\Code{Atom}(0 \leq x \leq 4) \;\judgfail
  \end{align*}
  The demonic assertion $\demonic\,\Code{Atom}(x = 1)$ fails because
  no superset of $r_x$ satisfies $\Code{Atom}(x = 1)$ exactly (since
  $r_x$ already contains $[x \mapsto 2]$, any superset violates
  exactness).  The angelic assertion $\angelic\,\Code{Atom}(0 \leq x
  \leq 4)$ fails because no sub-capability of $r_x$ satisfies
  $\Code{Atom}(0 \leq x \leq 4)$ exactly (that would require
  environments for $x = 0, 3, 4$ which $r_x$ does not contain).

  The disjointness distinction between $\chstar$ and $\chland$ also
  matters for these modalities.  Using the capabilities from
  Example~\ref{ex-capability-disjointness}:
  \begin{align*}
    &r_x \times r_y \models (\angelic\,\Code{Atom}(1 \leq x \leq 2)) \chstar (\angelic\,\Code{Atom}(3 \leq y \leq 4)) \;\judgok \\
    &r_{xy} \models (\angelic\,\Code{Atom}(1 \leq x \leq 2)) \chland (\angelic\,\Code{Atom}(3 \leq y \leq 4)) \;\judgok \\
    &r_{xy} \models (\angelic\,\Code{Atom}(1 \leq x \leq 2)) \chstar (\angelic\,\Code{Atom}(3 \leq y \leq 4)) \;\judgfail
  \end{align*}
  The product $r_x \times r_y$ satisfies the multiplicative
  conjunction because $x$ and $y$ are independently bound.  The
  dependent capability $r_{xy}$ satisfies the additive conjunction but
  not the multiplicative one, since its bindings cannot be decomposed
  into two disjoint sub-capabilities.
\end{example}

\paragraph{Sum Operator}
The sum $P \choplus Q$ holds at $r$ if $r$ can be split into two
sub-capabilities $r_1 + r_2 \sqsubseteq r$, one satisfying $P$ and
one satisfying $Q$.  This supports the context sum type $\tau_1
\ctoplus \tau_2$ introduced in \S\ref{sec:context-typing}, where a
single angelic context is partitioned across two control-flow branches.
The well-formedness requirement that capabilities are non-empty ensures
that neither branch is vacuous: every summand must witness at least one
environment.

\paragraph{Binding Reference Operator}
The binding reference operator $\forallM{x} P$ expresses how a
qualifier may depend on the specific binding of variable $x$ in the
current capability.  It holds at $r$ if, for every possible binding $v_x$
of $x$ in $r$, the sub-capability $\fiberOver{r}{[x \mapsto v_x]}$
satisfies $P$ with $x$ substituted by $v_x$.  This is needed because
$\Code{Atom}(\phi)$ has exact semantics: a qualifier $\phi$ mentioning
$x$ requires capability $r$ to contain \emph{all} environments
consistent with $x$'s full range, which conflicts with the constrained
range of $x$ expressed by other conjuncts.  The reference operator
resolves this by fixing the binding of $x$ before evaluating the
qualifier, enabling dependent qualifiers to refer to specific variable
bindings in the capability.

\begin{example}[Binding Reference]\label{ex-variable-reference}
  Consider the dependent angelic context $x{:}\urt{\Nat}{1 \leq \vnu
    \leq 2} \ctcomma y{:}\urt{\Nat}{0 \leq \vnu < x}$, where $y$'s
  range depends on $x$.  Its denotation is:
  \[
    \denot{x{:}\urt{\Nat}{1 \leq \vnu \leq 2} \ctcomma y{:}\urt{\Nat}{0 \leq \vnu < x}}
    = \angelic\,(1 \leq x \leq 2) \chland \forallM{x}\,\angelic\,(0 \leq y < x).
  \]
  The reference operator $\forallM{x}$ fixes the binding of $x$ before
  evaluating $\angelic\,(0 \leq y < x)$, so $y$'s required range is
  determined separately for each value of $x$.  Capabilities that
  model this formula include:
  \begin{align*}
    &r_1 \doteq \{[x \mapsto 1;\, y \mapsto 0]\} \\
    &r_2 \doteq \{[x \mapsto 2;\, y \mapsto 0];\, [x \mapsto 2;\, y \mapsto 1]\}
  \end{align*}
  and any sum $r_1 + r_2$ of compatible such capabilities.
\end{example}

\subsection{Type Denotation and Soundness}\label{subsec:denotation}

We now define the denotation of types and type contexts in terms of
the context logic, and state the soundness theorem.  We write
$\denot{\tau}_\Sigma : \Code{Term} \to \Code{Formula}$ for the type
denotation under basic type context $\Sigma$, where $r \models
\denot{\tau}_\Sigma\;e$ means that $e$ inhabits type $\tau$ under
capability $r$.  We omit $\Sigma$ when it is clear from context.

As in other refinement type systems, we require terms to be total and
well-typed in the basic type system.  We encode these side conditions
into the logic via the guard formula:
\begin{align*}
  P_{\Code{guard}} \doteq \forallM{\Dom{\Sigma}}\,\Code{Atom}(\Code{Total}(e) \land \Sigma \vdash_s e : \eraserf{\tau} \land \Sigma \vdash_{\Code{WF}} \tau).
\end{align*}
The reference operator $\forallM{\Dom{\Sigma}}$ is necessary here:
without it, $\Code{Atom}(\Sigma \vdash_s x : \Nat)$ would require the
capability to contain all possible bindings of $x$, making the basic
type judgement trivially true regardless of $x$'s actual range.  All
type denotations below implicitly include this guard, which we omit for
brevity.

The denotation of a term $e$ under capability $r$ captures all values
$e$ can reduce to across all environments in $r$:
\begin{align*}
  \denot{\mstep{e}{x}}_\Sigma &\doteq \forallM{\Dom{\Sigma}}\,\Code{Atom}(\mstep{e}{x}).
\end{align*}
The reference operator again fixes variable bindings before evaluating
the atomic reduction predicate, ensuring that in a nondeterministic
language the denotation correctly tracks which values are reachable
from which bindings.

\begin{definition}[Type Denotation]\label{def:type-denotation}
  \begin{align*}
    \denot{\ort{b}{\phi}}_\Sigma\;e &\doteq \denot{\mstep{e}{\vnu}}_\Sigma \chimpl \forallM{\Code{FV}(\ort{b}{\phi})}\,\demonic\,\phi \quad \text{($\vnu$ fresh)}
    \\\denot{\urt{b}{\phi}}_\Sigma\;e &\doteq \denot{\mstep{e}{\vnu}}_\Sigma \chimpl \forallM{\Code{FV}(\urt{b}{\phi})}\,\angelic\,\phi \quad \text{($\vnu$ fresh)}
    \\\denot{x{:}\tau_x \ctsarr \tau}_\Sigma\;e &\doteq \denot{\tau_x}_\Sigma\;x \chimpl \denot{\tau}_{\Sigma,x{:}\eraserf{\tau_x}}\;(e\;x) \quad \text{($x$ fresh)}
    \\\denot{x{:}\tau_x \ctwand \tau}_\Sigma\;e &\doteq \denot{\tau_x}_\Sigma\;x \chwand \denot{\tau}_{\Sigma,x{:}\eraserf{\tau_x}}\;(e\;x) \quad \text{($x$ fresh)}
    \\\denot{\tau_1 \ctinter \tau_2}_\Sigma\;e &\doteq \denot{\tau_1}_\Sigma\;e \chland \denot{\tau_2}_\Sigma\;e
    \\\denot{\tau_1 \ctunion \tau_2}_\Sigma\;e &\doteq \denot{\tau_1}_\Sigma\;e \chlor \denot{\tau_2}_\Sigma\;e
    \\\denot{\tau_1 \ctoplus \tau_2}_\Sigma\;e &\doteq \denot{\tau_1}_\Sigma\;e \choplus \denot{\tau_2}_\Sigma\;e
  \end{align*}
\end{definition}

The demonic and angelic type denotations follow directly from the
corresponding modalities: the result of evaluating $e$ (stored in
fresh variable $\vnu$) must satisfy $\phi$ under either the demonic or
angelic interpretation, with the reference operator fixing the bindings
of the free variables of $\phi$ before applying the modality.  The
function type denotations internalize the adjunction between function
types and context combinators: $\ctsarr$ uses additive implication
(the parameter binding is entangled with the closure context), while
$\ctwand$ uses multiplicative implication (the parameter binding is
disjoint from the closure context).

\begin{definition}[Type Context Denotation]\label{def:context-denotation}
  \begin{align*}
    \denot{\emptyset}_\Sigma &\doteq \chtrue
    \\\denot{x{:}\tau}_\Sigma &\doteq \denot{\tau}_\Sigma\;x
    \\\denot{\Gamma_1 \ctcomma \Gamma_2}_\Sigma &\doteq \denot{\Gamma_1}_\Sigma \chland \denot{\Gamma_2}_{(\Sigma,\,\eraserf{\Gamma_1})}
    \\\denot{\Gamma_1 \ctstar \Gamma_2}_\Sigma &\doteq \denot{\Gamma_1}_\Sigma \chstar \denot{\Gamma_2}_\Sigma
    \\\denot{\Gamma_1 \ctoplus \Gamma_2}_\Sigma &\doteq \denot{\Gamma_1}_\Sigma \choplus \denot{\Gamma_2}_\Sigma
  \end{align*}
\end{definition}

The dependent conjunction $\ctcomma$ maps to additive conjunction
$\chland$, reflecting that the second context may refer to bindings in
the first.  The disjoint conjunction $\ctstar$ maps to multiplicative
conjunction $\chstar$, requiring the two sub-capabilities to be
decomposable as a product — i.e., their variable bindings are
independent.

The semantic subtyping relations used in rules T-Sub and T-CtxSub are:
\begin{figure}[h]
  {\small
  \begin{flalign*}
    &\text{\textbf{Semantic subtyping}} && \fbox{$\Gamma \vdash \tau \subty \tau$ \quad $\Gamma \ctsubseteq{X} \Gamma$}
  \end{flalign*}
  \mprooftr{48mm}
  {$\forall e.\; \denot{\Gamma} \models \denot{\tau_1}\;e \chimpl \denot{\tau_2}\;e$}
  {Sub}
  {$\Gamma \vdash \tau_1 <: \tau_2$}
  \quad
  \mprooftr{66mm}
  {$\forall r.\; r \models \denot{\Gamma_1} \implies \exists r'.\; \projectTo{r}{X} = \projectTo{r'}{X} \land r' \models \denot{\Gamma_2}$}
  {CtxSub}
  {$\Gamma_1 \ctsubseteq{X} \Gamma_2$}
  }
  \caption{Semantic subtyping relations.}
  \label{fig:semantic-subtyping-relation}
\end{figure}

Type subtyping is encoded directly as logical implication in the
denotational semantics.  Context subtyping uses an explicit projection
$\projectTo{r}{X}$: the two capabilities agree on the variable set $X$
(the free variables of the term and type), but $r'$ may differ on
other variables, reflecting the fact that context subtyping only needs
to hold for the variables that actually appear in the judgement.

We conclude with the soundness results.

\begin{theorem}[Well-Formed Operator Context]\label{thm:well-formed-operator}
  Operator context $\Phi$ is well-formed iff for all primitive
  operators $\primop$:
  \begin{align*}
    \Phi(\primop) = \overline{x{:}\ort{b_x}{\phi_x}} \ctsarr \ort{b}{\phi} \ctinter \urt{b}{\phi}
    \quad\text{and}\quad
    \models \denot{\overline{x{:}\ort{b_x}{\phi_x}}} \chiff \denot{\ort{b}{\phi} \ctinter \urt{b}{\phi}}\;(\primop\;\overline{x}).
  \end{align*}
\end{theorem}

\begin{theorem}[Fundamental Theorem]\label{thm:fundamental}
  $\Gamma \vdash e : \tau$ implies $\denot{\Gamma}_\emptyset \models \denot{\tau}_{\eraserf{\Gamma}}\;e$.
\end{theorem}

\begin{theorem}[Denotational Soundness]\label{thm:soundness}
  $\emptyset \vdash e : \tau$ implies that for any fresh variable $x$,
  $\{[x \mapsto v] \mid \mstep{e}{v}\} \models \denot{\tau}\;x$.
\end{theorem}
